% Smutstitel, titelsida och kolofon.
%
% \bookversion är den release PDF:en byggdes för: `make VERSION=v2` skickar in den, och
% release-workflowet skickar in taggen. Utan den bär titelsidan bara byggdatumet.
\providecommand{\bookversion}{}
\newcommand{\bookdate}{\number\day\space\ifcase\month\or januari\or februari\or mars\or april\or
  maj\or juni\or juli\or augusti\or september\or oktober\or november\or december\fi\space
  \number\year}
\newcommand{\bookedition}{\ifx\bookversion\empty\else\bookversion\,·\,\fi\bookdate}

\thispagestyle{empty}
\vspace*{2.2in}
\begin{flushright}
  {\Large\bfseries Kommunikationsprotokoll och drivrutiner}
\end{flushright}
\cleardoublepage

\thispagestyle{empty}
\vspace*{1.1in}
\begin{flushleft}
  {\color{accent}\rule{0.9in}{2.5pt}}\par\vspace{1.4ex}
  {\fontsize{30}{36}\selectfont\bfseries Kommunikations-\\[0.15ex] protokoll och\\[0.15ex]
   drivrutiner\par}
  \vspace{2.2ex}
  {\Large\itshape Från eget protokoll till en testad\\CAN-drivrutin i C++\par}
  \vspace{1.0in}
  {\large Erik Pihl\par}
  \vfill
  {\small\color{muted} \bookedition\par}
\end{flushleft}
\newpage

\thispagestyle{empty}
\vspace*{\fill}
{\small
\noindent\textit{Kommunikationsprotokoll och drivrutiner}\\
Kursbok. Ett kapitel per föreläsning.

\medskip
\noindent Boken tar dig från ett självpåhittat protokoll, där varje mekanism skrivs för hand, till
en drivrutin mot en CAN-kontroller i ett FPGA som en parallellklass konstruerar samtidigt. Den
sista tredjedelen kopplar ihop de två halvorna på riktig hårdvara.

\medskip
\noindent Brödtexten är på svenska. All kod, alla kommentarer i koden, alla register- och
signalnamn är på engelska --- precis som i kursens övriga material och i koden du skriver.

\medskip
\noindent Text och figurer: Erik Pihl.

\medskip
\noindent Bokens text och figurer, liksom kursmaterialet den är satt från, är licensierade under
Creative Commons Erkännande 4.0 Internationell (CC BY 4.0). Du får kopiera, sprida och bearbeta
materialet i vilket syfte som helst, även kommersiellt, förutsatt att du anger upphovsperson och
licens. Licensen finns på \url{https://creativecommons.org/licenses/by/4.0/deed.sv}.
\medskip
\noindent Koden i bokens listningar får även användas under MIT-licensen, som är licensen för all
kod i kursens repository. Där finns också lektionsmaterialet och lösningsförslagen till övningarna.

\medskip
\noindent Koden är C++17. Drivrutinen håller sig till en delmängd som går att kompilera
friställt för en AVR32DB28, alltså utan heap, undantag och RTTI. AVR, AVR-Dx, Cyclone V och
DE0-CV är varumärken som tillhör sina respektive ägare.

\medskip
\noindent Satt med Lua\LaTeX\ i TeX Gyre Pagella och DejaVu Sans Mono.

\medskip
\noindent\textcolor{muted}{Den här utgåvan: \bookedition.}
\par}
\cleardoublepage
